%\ifodd\value{page}\hbox{}\newpage\fi
\chapter*{Ṛgveda 10.125 — Devī Sūkta, Mantra 1}
\addcontentsline{toc}{chapter}{Devī Sūkta 10.125.1}

\begin{sloka}
{\sanskritfont
\begin{tabular}{c}
ॐ अहं रुद्रेभिर्वसुभिश्चराम्यहमादित्यैरुत विश्वदेवैः ।\\
अहं मित्रावरुणोभा बिभर्म्यहमिन्द्राग्नी अहमश्विनोभा ॥१॥
\end{tabular}

\vspace{1.5em}

{\middlesize \iastfont
oṁ ahaṁ rudrebhir vasubhiś carāmy aham ādityair uta viśvadevaiḥ |\\
ahaṁ mitrāvaruṇobhā bibhar my aham indrāgnī aham aśvinobhā ||1||}

\vspace{1em}

{\middlesize
\anvaya{%
अहं रुद्रेभिः वसुभिः च रामि ।
अहं आदित्यैः उत विश्वदेवैः ।
अहं मित्रावरुणोभा बिभर्मि ।
अहं इन्द्राग्नी ।
अहं अश्विनोभा ।
}
}

\middlesize \iastfont \itmean{%
«Io mi muovo con i \textit{Rudra} e con i \textit{Vasu};  
io con gli \textit{Āditya} e anche con tutti gli Dei.  
Io sostengo entrambi, \textit{Mitra} e \textit{Varuṇa};  
io sono \textit{Indra} e \textit{Agni}, io sono entrambi gli \textit{Aśvin}.»%
}

}
\end{sloka}

\wordmeaning{%
\wordline{अहम्}
{\textit{aham}}
{ io; la voce che parla in prima persona, identificata con la potenza cosmica della Dea; }%

\wordline{रुद्रेभिः वसुभिः च रामि}
{\textit{rudrebhiḥ vasubhiḥ ca rām}}
{ mi muovo, agisco (\textit{rāmi}) con i \textit{Rudra} e con i \textit{Vasu}, due classi di divinità vediche; }%

\wordline{आदित्यैः उत विश्वदेवैः}
{\textit{ādityaiḥ uta viśvadevaiḥ}}
{ con gli \textit{Āditya} e anche (\textit{uta}) con tutti gli Dei (\textit{viśvadeva}); }%

\wordline{मित्रावरुणोभा बिभर्मि}
{\textit{mitrāvaruṇobhā bibharm}i}
{ sostengo (\textit{bibharmi}) entrambi (\textit{ubhā}) \textit{Mitra} e \textit{Varuṇa}, divinità dell’ordine e della legge cosmica; }%

\wordline{इन्द्राग्नी अश्विनोभा}
{\textit{indrāgnī aśvinobhā}}
{ io sono \textit{Indra} e \textit{Agni}; io sono entrambi gli \textit{Aśvin}, le coppie divine del vigore e del soccorso; }%
}

\textit{Devī Sūkta} 10.125 si apre con una dichiarazione radicale in prima persona. La voce che parla non è quella di una divinità particolare, ma di una potenza che si identifica con il funzionamento stesso del cosmo. L’uso reiterato di \textit{aham} non indica individualità, ma presenza onnipervasiva. La Dea afferma di “muoversi con” tutte le grandi classi divine, non come loro subordinata, ma come principio che le attraversa e le rende operative. Quando dice “io sostengo \textit{Mitra} e \textit{Varuṇa”}, la funzione è chiaramente cosmica: la legge, l’ordine e il patto sono portati e mantenuti dalla sua potenza. L’identificazione finale con \textit{Indra}, \textit{Agni} e con gli \textit{Aśvin} non è un’eliminazione delle differenze divine, ma la loro radice comune. La Dea non sostituisce gli dèi: li rende possibili. Questo primo mantra stabilisce così il tono dell’intero inno: la Dea come soggetto cosmico, voce dell’energia che precede e sostiene ogni funzione divina.

I \textit{Rudra} sono potenze tempestose e ambivalenti, legate alla forza distruttiva ma anche alla guarigione; i \textit{Vasu} incarnano invece gli aspetti stabili e luminosi della natura, come la luce, la terra e il fuoco benefico. Gli \textit{Āditya} sono divinità solari connesse alla legge cosmica e alla sovranità morale, mentre i \textit{Viśvadeva} rappresentano l’insieme indistinto di tutte le potenze divine. \textit{Mitra} e \textit{Varuṇa} presiedono all’ordine, al patto e alla giustizia cosmica; \textit{Indra} è il dio guerriero e sovrano, simbolo della forza che abbatte il caos; \textit{Agni} è il fuoco sacrificale, mediatore tra dèi e uomini; infine gli \textit{Aśvin}, divinità gemelle, sono associati al soccorso, alla guarigione e alla rapidità salvifica.